\section{Demonstration}
\label{sec:demo}

\begin{figure}[t]
  \centering
  \small
  \begin{tikzpicture}[
      node distance=4mm,
      stage/.style={draw, rounded corners=1pt, align=center,
                    minimum height=9mm, text width=24mm, inner sep=2pt},
      arrow/.style={->, line width=0.45pt}]
    \node[stage] (input) {Image\\+ regions};
    \node[stage, right=of input] (control) {Per-layer\\attention bias};
    \node[stage, right=of control] (passes) {Baseline +\\intervened passes};
    \node[stage, right=of passes] (output) {GPS results\\+ displacement};
    \draw[arrow] (input) -- (control);
    \draw[arrow] (control) -- (passes);
    \draw[arrow] (passes) -- (output);
  \end{tikzpicture}
  \caption{Interactive analysis flow. Regions may be drawn manually or loaded
  from WeDetect-Uni; both model outputs are presented together.}
  \label{fig:interaction_flow}
  \vspace{-2mm}
\end{figure}

Attendees operate the system directly through three short scenarios.

\subsubsection{Regional dependence.}
The attendee uploads a street-level image, marks a visual cue such as a sign or
licence plate, and suppresses its region. A large map displacement demonstrates
sensitivity to that evidence; a small displacement indicates that the current
prediction is stable under the intervention.

\subsubsection{Proposal comparison.}
For a benchmark example, cached WeDetect-Uni proposals can be applied
individually or as a union. The resulting displacements distinguish sensitivity
to one foreground object from sensitivity to the broader object-bearing region,
without a text prompt tailored to the image.

\subsubsection{Layer localisation.}
The attendee compares a single active layer with a multi-layer schedule.
Broader intervention often increases displacement without improving error,
mirroring Section~\ref{sec:experiments}. Resetting all biases to zero restores
the baseline exactly.

\subsection{Deployment and Relevance}

The live system stores GeoCLIP, its $10^5$-point GPS gallery, location
embeddings and cached proposals locally. The current interface obtains
basemap tiles and reverse-geocoded place names from public map services;
these provide cartographic context but are not part of model inference.

We bring the demonstration machine; it needs no GPU and runs inference
offline. The appendix states the setup requirements in full. The implementation and
experiment scripts will be made publicly available upon acceptance.
GeoProbe connects multimodal geolocation with causal explanation: rather than
presenting an attention map as evidence, it lets the attendee change an
attention pathway and observe the consequent geographic output. This
distinction is relevant to both multimedia verification and the privacy risks
of inferring location from untagged images.
